\hypertarget{appendix-87d7626c8d37}{}
\section*{Coordinate handedness and orientation}
\section{Coordinate handedness and orientation}\label{coordinate-handedness-and-orientation}

\textbf{Curriculum depth:} Introduced · \textbf{First seen:} Phase6 · \textbf{Latest development:} Phase6

\subsection{Definition and mechanism}\label{definition-and-mechanism}

Handedness is the sign of an ordered basis relative to a chosen orientation. A reflection reverses it; a proper rotation preserves it. It does not specify up axis, camera direction or matrix storage order.

The sign of \((\mathbf x\times\mathbf y)\cdot\mathbf z\) classifies a nondegenerate frame. Swapping two axes flips the sign.

\subsection{Prerequisites}\label{prerequisites}

\begin{itemize}
\tightlist
\item
  \hyperref[appendix-2391db8052b1]{Cross products and oriented area} --- Oriented area supplies face normals, frame signs and geometric determinants.
\end{itemize}

\subsection{Used by}\label{used-by}

Later curriculum sources can deepen this same node; no future coverage is marked complete.

\subsection{Recall and next depth}\label{recall-and-next-depth}

Trace winding through a mirror while preserving ordinary distances. Full production proofs, target-specific integration and learner mastery require further evidence.

\subsection{Source and navigation}\label{source-and-navigation}

\href{Continuous\%20Notes/Computing\%20Foundations\%20-\%20From\%20Bits\%20to\%20C++\%20and\%20Systems\%20APIs\#Oriented\%20area\%20tells\%20us\%20which\%20side\%20of\%20the\%20panel\%20is\%20front}{Teaching narrative} · \href{}{Coverage and source gap (01-Foundations-I, appendix-fd090d982a15 § Phase 6 - Parts 501-600)} · \href{}{Evidence and limits (01-Foundations-I, appendix-a103e7846d21 § Phase 6 - Source analysis and verification)} · \hyperref[appendix-4ad47bdd269a]{Dependency Map}
